\section{Conclusion}
\label{sec:conclusion}
In this work, we introduced the Shadow Enhanced Greedy Quantum Eigensolver (SEGQE), a greedy, shadow-assisted hybrid algorithm for measurement-efficient ground-state preparation. 
The central feature of SEGQE is its favorable measurement scaling. At each iteration, a single batch of classical shadows suffices to evaluate a large number of candidate gates entirely in classical post-processing. We proved rigorous worst-case bounds showing that the per-iteration measurement cost depends only logarithmically on the number of candidate gates. In addition, our numerical experiments on transverse-field Ising models and ensembles of random local Hamiltonians indicate that the required number of iterations grows approximately linearly with system size. Together, these results suggest an overall scaling behavior that is particularly well aligned with early fault-tolerant regimes, where logical measurement rates are relatively slow while moderately deep logical circuits are feasible.
Furthermore, SEGQE cleanly separates quantum data acquisition from the computationally intensive search over gate additions, enabling the efficient integration of high-performance computing resources and shifting the dominant optimization cost to classical hardware. 

The flexibility of the SEGQE framework further allows straightforward adaptation to hardware-native logical gate sets and fault-tolerant compilation constraints, making it compatible with realistic fault-tolerant architectures. Future work includes tightening instance-dependent performance guarantees, extending the framework to related tasks such as excited-state preparation, and developing parallelizable classical optimization strategies tailored to shadow-based objective estimation.
Overall, SEGQE illustrates how classical shadows can serve as a core primitive in adaptive, measurement-efficient state-preparation strategies for fault-tolerant quantum algorithms.
